\par
\noindent\hspace*{\parindent}\textls[100]{5. Existenzbeweis für Schnittpunkte von \(U=0\) und \(T=0\).}
Um nun die Existenz von Schnittpunkten von \(U=0\) und \(T=0\) zu beweisen,
definieren wir auf jedem Elementarbogen von \(U=0\) eine Fortschreitungsrichtung
durch die Festsetzung, dass beim Fortschreiten längs eines
solchen Bogens das Gebiet mit \(U^{(1)}>0\) zur rechten Hand bleiben soll. — Dies
ist möglich, da ja jeder Elementarbogen von \(U=0\) ein Gebiet mit \(U^{(1)}>0\)\par
\clearpage\SourceLeaf{325}\ArithmeticHead{16}{L}{A. OSTROWSKI, ÜBER DEN ERSTEN UND VIERTEN GAUSSSCHEN BEWEIS USW.}
\noindent von einem Gebiet mit \(U^{(1)}<0\) trennt. Gelangen wir aber so zu einem Punkte,
in dem mehrere Elementarbogen zusammentreffen, so sei festgesetzt, dass wir
von jedem Elementarbogen auf den \textls[100]{nächsten von rechts} hinübergehen.
Dabei hat man stets ein Gebiet mit \(U^{(1)}>0\) zur rechten Hand. Treten wir
nun in irgend einem Intervall \((4k+1)\) längs der Kurve \(U=0\) in den Kreis
\((R)\) ein, so werden wir nach unserer Festsetzung des Fortschreitungssinnes keinen
Elementarbogen zweimal durchlaufen können und werden daher in irgend
einem Intervall \((4k+3)\) aus dem Kreise austreten. Wir erhalten daher im
ganzen \(n\) sich ins Unendliche erstreckende Zweige der Kurve \(U=0\), die in
den Intervallen \((4k+1)\) in den Kreis eintreten und in den Intervallen \((4k+3)\)
aus ihm austreten. Wir wollen einen Zweig, der im Intervalle \((4k+1)\) in
den Kreis eintritt, durch \([k]\) bezeichnen. Nach unseren Festsetzungen über
den Fortschreitungssinn können zwei von diesen Zweigen nie einen Elementarbogen
gemeinsam haben, sondern höchstens einzelne Punkte. (Aber auch in
solchen Punkten \textls[100]{durchsetzen} sich, wie aus unseren Festsetzungen folgt, solche
Zweige nie, sondern stossen nur aneinander an.)
\par
Durch das Vorhergehende ist es noch nicht ausgeschlossen, dass es Elementarbogen
gibt, die keinem der sich ins Unendliche erstreckenden Zweige
angehören, Wir werden bald zeigen, dass es solche Bogen nicht geben kann.
Für unseren Zweck ist dies aber sogar unnötig. Denn in dem Punkte im
Intervall \((4k+1)\), in dem ein Zweig \([k]\) in den Kreis eintritt, ist das Vorzeichen
von \(T\) positiv, im Austrittspunkte in einem Intervalle \((4k'+3)\) aber
negativ. Daher verschwindet \(T\) auf dem Zweige \([k]\) von \(U=0\) in wenigstens
einem Punkte, und damit ist gezeigt, dass auf jedem der \(n\) Zweige \([k]\) wenigstens
eine Wurzel von \(f(z)=0\) liegt. Es ist sogar die Existenz von \(n\) Wurzeln
von \(f(z)\) nachgewiesen, allerdings unter der Annahme, dass die \(n\) Zweige \([k]\)
keinen Punkt gemeinsam haben. Um auch für den Fall, dass gewisse Zweige
\([k]\) aneinander stossen, die Existenz von genau \(n\) Wurzeln von \(f(z)\) nachzuweisen,
hat man nur zu zeigen, dass, wenn in einem Punkte, in dem \(l\) Zweige
von \(U=0\) aneinander stossen, wenigstens eine Wurzel von \(f(z)\) liegt, in ihm
genau \(l\) Wurzeln von \(f(z)\) liegen. Doch wollen wir darauf nicht näher eingehen\footnote[3]{Vergl. Werke III, S. 83—84.}.
\par\clearpage\SourceLeaf{326}\ArithmeticHead{17}{R}{A. OSTROWSKI, ÜBER DEN ERSTEN UND VIERTEN GAUSSSCHEN BEWEIS USW.}
\HistoricalFoot{X\textsubscript{2} Abh. 3.}{3}
\noindent\hspace*{\parindent}\textls[100]{6. Weitere Eigenschaften der Kurve \(U=0\).} Wir wollen noch
der Vollständigkeit wegen zeigen, dass jeder Elementarbogen von \(U=0\) zu
einem der Zweige \([k]\) gehört. Denn gehörte ein Elementarbogen \(E\) keinem
der \(n\) Zweige \([k]\) an, so könnten wir eine solche reelle Konstante \(j\) finden,
dass \(f(z)+j\) auf \(E\) eine Nullstelle hätte, dagegen in keinem der Punkte, in
denen die Zweige \([k]\) von \(U=0\) eventuell aneinander stossen, verschwände.
Da die \(U\)-Kurve von \(f(z)+j\) mit der von \(f(z)\) übereinstimmt, so hätte \(f(z)+j\)
auf jedem der \(n\) Zweige \([k]\) eine Wurzel, also im ganzen \(n\) Wurzeln, ausserdem
aber noch eine Wurzel auf \(E\), also im ganzen wenigstens \(n+1\) Wurzeln,
\(f(z)\) ist aber vom Grade \(n\).
\par
Aus der damit bewiesenen Tatsache können wir in einfacher Weise eine
weitere Tatsache folgern, die den Verlauf der Kurve \(U=0\) im Sinne der
Analysis Situs charakterisiert. Es ist nämlich unmöglich, aus den Elementarbogen
von \(U=0\) einen geschlossenen Zug zu bilden. Die Kurve \(U=0\)
zerfällt also, topologisch gesprochen, in eine Anzahl von \textls[100]{Bäumen}, die miteinander
keinen Punkt gemeinsam haben. Dies lässt sich sofort aus der Potentialeigenschaft
von \(U\) erschliessen, da eine reguläre Potentialfunktion, die
auf der ganzen Berandung eines ganz im Endlichen liegenden abgeschlossenen
Bereiches gleich \(0\) ist, in der ganzen Ebene verschwinden muss. Es lässt sich
dies aber auch elementar wie folgt zeigen:
\par
Wäre es möglich, aus den Elementarbogen von \(U=0\) einen geschlossenen
Zug zu bilden, so könnte man, wie man leicht einsieht, auch einen solchen
doppelpunktlosen geschlossenen Zug \(E_1E_2\ldots E_s\) bilden, in dessen Innerem
kein Elementarbogen von \(U=0\) verläuft. Wir wollen annehmen, dass beim
Durchlaufen von \(E_1E_2\ldots E_s\) dessen Inneres zur rechten Hand bleibt, was
durch Umnumerieren der Bogen \(E_1\ldots E_s\) jedenfalls zu erreichen ist. Ist
nun im Inneren von \(E_1\ldots E_s\) \(U^{(1)}\geqq0\), und schreiten wir längs eines Bogens
\(E_1\) gemäss dem früher festgesetzten Fortschreitungssinne fort, so werden wir nach
diesen Festsetzungen stets auf dem Zuge \(E_1\ldots E_s\) bleiben, so dass \(E_1\) keinem
der \(n\) Züge \([k]\) angehört. Dies widerspricht aber der oben bewiesenen Tatsache.
— Ist aber im Innern von \(E_1\ldots E_s\) \(U^{(1)}\leqq0\), so ändern wir unsere
Festsetzung des Fortschreitungssinnes in 5. dahin ab, dass beim Fortschreiten
\(U^{(1)}\) zur rechten Hand negativ sein soll. Dann bleibt unsere Überlegung mutatis
mutandis in Kraft. —
\par\clearpage\SourceLeaf{327}\ArithmeticHead{18}{L}{A. OSTROWSKI, ÜBER DEN ERSTEN UND VIERTEN GAUSSSCHEN BEWEIS USW.}
Wir haben oben anstatt \(U^{(1)}>0,\ U^{(1)}<0\) geschrieben \(U^{(1)}\geqq0,\ U^{(1)}\leqq0\),
um dem Vorkommen von isolierten Nullstellen von \(U^{(1)}\) Rechnung zu tragen. Es
ist aber leicht nach derselben Methode zu zeigen, dass es isolierte Nullstellen
von \(U\) nicht gibt. Denn wäre \((x_0,y_0)\) eine solche, so gäbe es eine reelle
Konstante \(j\) der Art, dass \(f(z)+j\) in \((x_0,y_0)\) eine Wurzel hätte. Ändern wir
die \textls[100]{reelle} Konstante \(j\) beliebig wenig ab, so folgte aus der Stetigkeit der
Wurzeln als Funktionen der Koeffizienten, dass die Nullstelle \((x_0,y_0)\) in eine
andere übergeht, die von \((x_0,y_0)\) beliebig wenig entfernt bliebe. Daher müsste
\(U\) in beliebiger Nähe von \((x_0,y_0)\) Nullstellen haben, so dass \((x_0,y_0)\) nicht
isoliert sein kann.
\par\bigskip
\begin{center}\rule{3.4cm}{.4pt}\end{center}
\vspace{1.2cm}
\begin{center}\textbf{Berichtigungen zu der Abh. 3 (Gauss' Werke Bd. X, 2).}\end{center}
{\small
\begin{tabularx}{\linewidth}{@{}r r X@{}}
Seite&6,&Zeile 8 v. u. ist nach Veränderlichen einzufügen: »keinen von nur einer Veränderlichen abhängigen Faktor und«\\
»&11,&Zeile 4 v. u. ist nach \(x-x'\) einzufügen: »oder \(y-y'\)«\\
»&14,&Zeile 6 v. u. ist nach \(x=x'\) einzufügen: »und \(y=y'\)«\\
»&15,&Zeile 12 v. u. ist nach \(x=x_1\) einzufügen: »oder \(y=w\)«
\end{tabularx}\par}
\vfill
\begin{center}\small Göttingen, Druck der Dieterichschen Universitäts-Buchdruckerei (W. Fr. \textls[100]{Kaestner}).\end{center}
\par
